\section{Discussion}
\label{sec:discussion}

\paragraph{Where analytic ends and trained begins.} Frobenius reconstruction of a single weight matrix is capped by Eckart-Young. Every trained rank-$k$ approximator we compared against loses or ties on 24 / 24 modules; no amount of training can push past the analytic ceiling. The moment the objective breaks Eckart-Young ---nonlinear downstream composition, multi-D concept alignment measured against ground truth, causal preservation across layers --- training \emph{can} help. Our non-Frobenius extension formalizes this and demonstrates 73.9\% mean MSE reduction on adversarial constructions. Practitioners of interpretability should therefore use analytic \svdomp\ as the baseline; use training only when the objective explicitly requires composition or downstream signal.

\paragraph{Substrate matters.} Our results are on \emph{weight} decomposition. \bsf's original result is on \emph{activation} decomposition of a vision model. Activation concepts are known to be non-aligned with principal directions of variance (this is why SAEs beat PCA in practice), so we do not claim that \svdomp\ would beat \bsf\ on activations. We do claim \svdomp\ dominates trained methods on weight decomposition; and we claim the stable-rank plateau BSF interprets as "concepts are 2--4 dimensional" is a property of activations rather than of the featurizer, because our analytic blocks reproduce the plateau exactly.

\paragraph{Limitations.} (i) \svdomp\ requires the weight matrix; it does not apply to methods that only expose activations. (ii) On the six attention $v_\text{proj}$ / $o_\text{proj}$ modules where singular spectra are compressed, Davis-Kahan predicts and we observe reduced stability; \vpd\ wins on those. (iii) All experiments in Sections \ref{sec:vs-vpd}--\ref{sec:non-frobenius-exp} run on Goodfire's 67M \textsc{LlamaSimpleMLP}. Whether the Pareto dominance transfers cleanly to modern reasoning-model scales (7B+, chain-of-thought-trained) requires verification against those models; the efficiency argument (Sec.\ \ref{sec:efficiency}) suggests it is precisely at large scale that \svdomp's zero training cost becomes materially valuable.

\paragraph{Conclusion.} For efficient interpretability of large-model weights, the analytic SVD dictionary supplies the Frobenius-optimal decomposition at zero training cost. Trained probes recover only what the SVD already provides; where they genuinely help is on non-Frobenius composed objectives. Our results reproduce and generalize \bsf's stable-rank finding without training, position analytic decomposition as the correct baseline for interpretability at scale, and provide a small trainable correction for cases where the objective exceeds the Eckart-Young cap.
